\documentclass[12pt,a4paper]{article}

% Niezbędne pakiety
\usepackage[T1]{fontenc}
\usepackage[utf8]{inputenc}
\usepackage[polish]{babel}
\usepackage{geometry}
\usepackage{graphicx}
\usepackage{tabularx}
\usepackage{xcolor}
\usepackage{setspace}
\usepackage{hyperref}

% Ustawienia marginesów i interlinii (dla lepszej czytelności i objętości)
\geometry{margin=2.5cm}
\onehalfspacing

% Konfiguracja linków
\hypersetup{
    colorlinks=true,
    linkcolor=black,
    filecolor=magenta,      
    urlcolor=blue,
    pdftitle={Dokumentacja Funkcjonalna - Java Swing},
}

\begin{document}

% ==================== STRONA TYTUŁOWA ====================
\begin{titlepage}
    \centering
    \vspace*{3cm}
    
    {\scshape\LARGE Projekt JIMP 2026 \par}
    \vspace{1.5cm}
    {\huge\bfseries Dokumentacja Funkcjonalna\par}
    \vspace{0.5cm}
    {\Large Aplikacja z Graficznym Interfejsem Użytkownika (Java Swing) \par}
    
    \vspace{2cm}
    \textbf{Moduł analityczny i wizualizacja podziału grafów planarnych}
    
    \vfill
    
    \begin{flushright}
        \large
        \textbf{Autorzy:}\\
        Ksawey Kuczyński\\
        Olaf Wodecki
    \end{flushright}
    

\end{titlepage}

\newpage
% ==================== SPIS TREŚCI ====================
\tableofcontents
\newpage

% ==================== WSTĘP ====================
\section{Wstęp i Cel projektu}
Niniejszy dokument stanowi specyfikację funkcjonalną drugiej części projektu realizowanego w ramach przedmiotu JIMP. Aplikacja została stworzona w języku Java z wykorzystaniem biblioteki graficznej Swing. 

Głównym celem programu jest dostarczenie użytkownikowi końcowemu intuicyjnego narzędzia okienkowego do wczytywania, analizy oraz wizualizacji grafów. Aplikacja ściśle współpracuje z modułem obliczeniowym napisanym w języku C, potrafiąc wczytywać wygenerowane przez niego pliki wynikowe (zarówno w formacie tekstowym, jak i zoptymalizowanym formacie binarnym).

Kluczową funkcjonalnością aplikacji jest \textbf{dokonywanie podziału grafu} na zadaną przez użytkownika liczbę rozłącznych części, przy zachowaniu równomierności podziału określanej za pomocą procentowego marginesu błędu. Wyniki tych operacji są natychmiastowo odzwierciedlane na ekranie, gdzie poszczególne podgrafy zyskują unikalne kolory.

\section{Wymagania systemowe i uruchomienie}
Aplikacja została zaprojektowana jako rozwiązanie wieloplatformowe (\textit{cross-platform}), co wynika z natury środowiska uruchomieniowego Java.
\begin{itemize}
    \item \textbf{Środowisko uruchomieniowe:} Java Runtime Environment (JRE) w wersji 11 lub nowszej.
    \item \textbf{Interfejs graficzny:} Środowisko graficzne (X11/Wayland w systemach Linux, Windows DWM lub macOS WindowServer).
    \item \textbf{Uruchomienie:} Program nie wymaga instalacji. Uruchamiany jest z poziomu wiersza poleceń wywołaniem: \texttt{java -jar GraphPartitioner.jar} lub poprzez dwukrotne kliknięcie pliku wykonywalnego JAR.
\end{itemize}

\newpage
% ==================== ARCHITEKTURA GUI ====================
\section{Architektura Interfejsu Użytkownika}
Interfejs graficzny został zaprojektowany z myślą o minimalizmie i ergonomii. Okno główne programu (reprezentowane przez klasę \texttt{JFrame}) zostało logicznie podzielone na trzy niezależne obszary operacyjne.

\subsection{Pasek Menu (Menu Bar)}
Znajduje się w górnej krawędzi okna aplikacji. Odpowiada za operacje wejścia/wyjścia na plikach.
\begin{itemize}
    \item \textbf{Plik $\rightarrow$ Wczytaj graf (TXT):} Otwiera okno wyboru pliku (\texttt{JFileChooser}) z filtrem rozszerzeń \texttt{.txt} lub \texttt{.csrrg}.
    \item \textbf{Plik $\rightarrow$ Wczytaj graf (BIN):} Pozwala na szybkie załadowanie wielkich grafów zrzucanych binarnie przez moduł w C.
    \item \textbf{Plik $\rightarrow$ Zapisz wynik podziału:} Eksportuje pokolorowany i podzielony graf z powrotem do pliku tekstowego na dysku.
    \item \textbf{Zakończ:} Bezpieczne zamknięcie aplikacji.
\end{itemize}

\subsection{Panel Narzędziowy (Control Panel)}
Znajduje się z lewej strony okna głównego i zawiera interaktywne kontrolki formularza. To tutaj użytkownik decyduje o parametrach działania algorytmu.

\begin{table}[h]
    \centering
    \begin{tabularx}{\textwidth}{|l|X|}
    \hline
    \textbf{Kontrolka} & \textbf{Opis i zachowanie} \\ \hline
    Pole tekstowe (Partycje) & Służy do wpisania docelowej liczby części, na które ma zostać podzielony graf. Wartość domyślna to 2. \\ \hline
    Suwak (Margines \%) & Pozwala zdefiniować procentowy margines tolerancji wielkości podziałów (od 0\% do 100\%). Wartość domyślna to 10\%. \\ \hline
    Checkbox (Wagi) & Przełącznik ukrywający lub pokazujący etykiety wag na krawędziach grafu w celu poprawy czytelności. \\ \hline
    Przycisk \textit{Podziel Graf} & Główny przycisk wyzwalający algorytm. Po jego kliknięciu program waliduje dane i przerysowuje płótno. \\ \hline
    \end{tabularx}
    \caption{Opis kontrolek panelu narzędziowego}
\end{table}

\subsection{Przestrzeń Robocza (Graph Canvas)}
Najważniejszy element aplikacji, zajmujący centralną i największą część okna. Jest to dedykowany komponent, w którym nadpisano metodę \texttt{paintComponent(Graphics g)}. 

Dzięki samodzielnej implementacji renderowania, przestrzeń robocza pozwala na:
\begin{itemize}
    \item Skalowanie i wyśrodkowywanie grafu niezależnie od rozmiarów okna.
    \item Przypisywanie unikalnych kolorów wierzchołkom należącym do różnych partycji (np. zbiór A = czerwony, zbiór B = niebieski, zbiór C = zielony).
    \item Rysowanie krawędzi łączących odpowiednie wierzchołki.
\end{itemize}

\newpage
% ==================== SCENARIUSZE UŻYCIA ====================
\section{Scenariusze użycia i przykłady wizualne}

Poniżej przedstawiono standardowy proces pracy z programem z perspektywy końcowego użytkownika.

\subsection{Krok 1: Wczytanie surowych danych}
Użytkownik po uruchomieniu programu wybiera z paska menu opcję wczytania danych. Po załadowaniu pliku wygenerowanego przez aplikację C, na ekranie pojawia się struktura grafu bazowego. W tym momencie wszystkie wierzchołki mają jednakowy kolor, ponieważ podział nie został jeszcze dokonany.

% MIEJSCE NA ZRZUT EKRANU 1
\begin{figure}[h]
    \centering
    % \includegraphics[width=0.8\textwidth]{screen1.png}
    % Zastępcze pole:
    \textcolor{gray}{\rule{0.8\textwidth}{6cm}}
    \caption{Widok głównego okna po wczytaniu bazowego grafu (przed podziałem).}
\end{figure}

\subsection{Krok 2: Parametryzacja i Podział}
W panelu narzędziowym użytkownik ustala liczbę części (np. 3) oraz margines (np. 15\%). Po wciśnięciu przycisku \textit{,,Podziel Graf''}, aplikacja przetwarza strukturę. Rysunek aktualizuje się, a wierzchołki przybierają 3 różne kolory oznaczające ich przynależność.

% MIEJSCE NA ZRZUT EKRANU 2
\begin{figure}[h]
    \centering
    % \includegraphics[width=0.8\textwidth]{screen2.png}
    % Zastępcze pole:
    \textcolor{gray}{\rule{0.8\textwidth}{6cm}}
    \caption{Wizualizacja grafu podzielonego na 3 rozłączne części (reprezentowane różnymi kolorami).}
\end{figure}

\newpage
% ==================== BŁĘDY I ZABEZPIECZENIA ====================
\section{Sytuacje wyjątkowe i komunikaty o błędach}

Aby zapewnić stabilność środowiska i zapobiec awariom wynikającym z pomyłek użytkownika, aplikacja na bieżąco weryfikuje stany komponentów. W przypadku wykrycia nieprawidłowości, interfejs graficzny wyświetla wyskakujące okna dialogowe (\textit{JOptionPane}) blokujące działanie do czasu usunięcia problemu.

Poniższa tabela zbiera wszystkie kluczowe wyjątki obsługiwane przez program:

\vspace{0.5cm}
\begin{table}[h]
    \centering
    \renewcommand{\arraystretch}{1.5}
    \begin{tabularx}{\textwidth}{|p{3.5cm}|X|p{3.5cm}|}
    \hline
    \textbf{Sytuacja} & \textbf{Opis zachowania programu i przyczyna} & \textbf{Reakcja GUI} \\ \hline
    
    \textbf{Błąd odczytu I/O} & Użytkownik wskazał plik, który nie istnieje, został usunięty w trakcie odczytu lub aplikacja nie ma uprawnień do jego otwarcia. & Okno błędu: \textit{,,Błąd systemu: Nie można otworzyć wskazanego pliku.''} Płótno pozostaje puste. \\ \hline
    
    \textbf{Zły format pliku} & Zawartość wybranego pliku nie zgadza się ze specyfikacją (np. podano plik MP3 zamiast pliku z grafem, lub plik tekstowy ma uszkodzone linie). & Okno błędu: \textit{,,Zły format danych. Upewnij się, że plik pochodzi z aplikacji C.''} \\ \hline
    
    \textbf{Niewłaściwe parametry wejścia} & Użytkownik wpisał w polu ,,Liczba części'' wartość ujemną, wartość $1$ lub tekst (np. litery). & Okno ostrzegawcze: \textit{,,Podano niewłaściwe parametry. Liczba części musi być liczbą całkowitą $>1$.''} Przycisk podziału jest ignorowany. \\ \hline
    
    \textbf{Zła kolejność operacji} & Użytkownik próbuje wcisnąć przycisk ,,Podziel Graf'' lub zmienić widok zanim jakikolwiek plik zostanie wczytany do pamięci. & Komunikat: \textit{,,Najpierw załaduj plik z grafem używając paska Menu.''} \\ \hline
    
    \end{tabularx}
    \caption{Zestawienie sytuacji wyjątkowych i reakcji interfejsu graficznego}
\end{table}

\section{Podsumowanie funkcjonalności}
Zaprojektowany interfejs graficzny ściśle realizuje wszystkie założenia specyfikacji. Dzięki wyizolowaniu paska poleceń, panelu ustawień i własnego komponentu renderującego, aplikacja jest wysoce skalowalna. Daje to w przyszłości możliwość łatwego dodawania kolejnych opcji (np. powiększania fragmentów grafu rolką myszy) bez konieczności przebudowywania układu okna głównego.

\end{document}